\section{Smoothed sums and their Mellin transforms}
\label{sec:mellin}

\begin{lemma}\label{lem:mellin}
For $\Re s>1$ one has
\begin{align}
 \int_1^\infty G_\chi(x)x^{-s-1}\,dx
  &=-\frac{L'(s,\chi)/L(s,\chi)
                  -sL'(1,\chi)/L(1,\chi)}{s(s-1)}
                       &&(\chi\ne1),\label{eq:mellin-chi}\\
 \int_1^\infty G_1(x)x^{-s-1}\,dx
  &=-\frac{\zeta'(s)/\zeta(s)}{s(s-1)}
      -\frac1{(s-1)^2}+\frac{1+\gamma}{s-1}.
                       \label{eq:mellin-zeta}
\end{align}
The expressions on the right extend meromorphically to $\Re s>0$,
are holomorphic at $s=1$, and have, at a zero $\rho$ of the
corresponding $L$-function in $0<\Re\rho<1$, residue
\begin{equation}\label{eq:residue}
                    -\frac{m_\chi(\rho)}{\rho(\rho-1)},
\end{equation}
where $m_\chi(\rho)$ is the positive integer multiplicity of that zero.
All their poles in $\Re s>0$ are of this form and are simple.
\end{lemma}

\begin{proof}
For $n\ge1$ and $\sigma>1$, direct integration gives
\[
 \int_1^\infty(x/n-1)_+x^{-s-1}\,dx
 =\frac{n^{-s}}{s(s-1)}.
\]
The same calculation with $\sigma$ in place of $s$ bounds the integral
of the absolute value by $n^{-\sigma}/(\sigma(\sigma-1))$.
Consequently the exchange of sum and integral is justified by
$\sum\Lambda(n)n^{-\sigma}<\infty$. Logarithmic differentiation of
the absolutely convergent Euler product gives
\[
       \sum_{n\ge1}\Lambda(n)\chi(n)n^{-s}
                         =-\frac{L'(s,\chi)}{L(s,\chi)}.
\]
Finally,
$\int_1^\infty x^{-s}\,dx=(s-1)^{-1}$ and
$\int_1^\infty x^{-s}\log x\,dx=(s-1)^{-2}$.
These identities prove both formulas, including absolute convergence
of the integrals for $\sigma>1$.

In \eqref{eq:mellin-chi} the numerator vanishes at $s=1$; its
apparent singularity there is removable. For \eqref{eq:mellin-zeta},
use $\zeta'/\zeta=-1/(s-1)+\gamma+O(s-1)$.
The coefficients of both $(s-1)^{-2}$ and $(s-1)^{-1}$ cancel.
If $L(s,\chi)=(s-\rho)^m h(s)$ with $h$ holomorphic and
$h(\rho)\ne0$, then $L'/L=m/(s-\rho)+h'/h$.
Since $\rho(\rho-1)\ne0$, this proves \eqref{eq:residue}.
There are no other singularities in the indicated half-plane, by
the classical analytic properties stated above. The zeros of a
nonzero holomorphic function are isolated, so the pole set is locally
finite.
\end{proof}


\begin{lemma}\label{lem:positive-family}
Let $\chi_1,\ldots,\chi_r$ be primitive characters, possibly including
the conductor-one character, and let $c_j>0$ be real numbers.
Set $H(x)=\sum_jc_j\Re G_{\chi_j}(x)$ and $g(u)=H(e^u)$.
Then $|g(u)|\ll e^{2u}$ for $u\ge0$. Its Laplace transform extends
meromorphically to $\Re s>0$. At every zero $\rho$ of one of the
$L(s,\chi_j)$ in that half-plane the transform has a genuine simple
pole, with residue
\begin{equation}\label{eq:family-residue}
 -\frac{1}{\rho(\rho-1)}
    \sum_{j=1}^r\frac{c_j}{2}
                 \bigl(m_{\chi_j}(\rho)+m_{\bar\chi_j}(\rho)\bigr).
\end{equation}
Here a multiplicity is zero if the corresponding function does not
vanish. The residue at a real pole is positive.
\end{lemma}

\begin{proof}
The bound $\Lambda(n)\le\log n$ gives
$|S_\chi(x)|\le x\sum_{2\le n\le x}(\log n)/n
\ll x(1+\log^2x)\ll x^2$. The correction terms in
\eqref{eq:Gchi}--\eqref{eq:Gone} obey the same bound for each fixed
character. Thus $g$ is continuous and has the stated growth.

Conjugating the absolutely convergent Dirichlet series, and then
continuing analytically, gives
$L(s,\bar\chi)=\overline{L(\bar s,\chi)}$ and
$G_{\bar\chi}(x)=\overline{G_\chi(x)}$.
Hence $\Re G_\chi=(G_\chi+G_{\bar\chi})/2$.
The substitution $x=e^u$ identifies each Mellin integral in
Lemma~\ref{lem:mellin} with the corresponding Laplace integral.
Formula~\eqref{eq:family-residue} follows by addition. Its integer
multiplicities are nonnegative, and at least one is positive at a
zero under consideration. Thus the residue cannot cancel.
For real $0<\rho<1$, the factor $-1/(\rho(\rho-1))$ is positive.
\end{proof}

